\item \textbf{Problema de repaso.}
Un objeto de $12.0\text{ kg}$ cuelga en equilibrio de una cuerda con una longitud total de $L = 5.00\text{ m}$ y una densidad de masa lineal $\mu = 0.00100\text{ kg/m}$. La cuerda se enrolla alrededor de dos poleas ligeras sin fricción separadas una distancia $d = 2.00\text{ m}$ (figura P17.47a).
\begin{enumerate}
    \item Determine la tensión en la cuerda.
    \item ¿A qué frecuencia debe vibrar la cuerda entre las poleas para formar el patrón de onda estacionaria que se muestra en la figura P17.47b?
\end{enumerate}